% TOP_PROBLEM: {"problemId": "problem.murres-chow-kunneth-conjectures", "canonicalTitle": "Murre's Chow-Kunneth Conjectures", "releaseRank": 153, "primaryDomain": "number_theory_arithmetic_geometry", "primaryDomainLabel": "Number theory & arithmetic geometry", "displayStatus": "open_with_solved_subcases", "releaseStatus": "open_with_solved_subcases"}
% !TeX program = lualatex
% Definitions and sources retained from research_batch_151-153.tex; research is in GPT_6_Astra_Ultra.
\documentclass[11pt]{article}
\usepackage[margin=25mm]{geometry}
\usepackage{fontspec}
\setmainfont{Latin Modern Roman}
\usepackage{amsmath,amssymb,mathtools}
\usepackage{unicode-math}
\setmathfont{Latin Modern Math}
\usepackage{enumitem,needspace}
\usepackage{xurl,hyperref}
\hypersetup{colorlinks=true,urlcolor=blue}
\setcounter{secnumdepth}{1}
\setlength{\parindent}{0pt}
\setlength{\parskip}{5pt}
\setlength{\emergencystretch}{3em}
\setlist[itemize]{leftmargin=3em,itemsep=5pt,topsep=4pt,parsep=0pt}
\allowdisplaybreaks
\newcommand{\N}{\mathbb{N}}
\newcommand{\Z}{\mathbb{Z}}
\newcommand{\Q}{\mathbb{Q}}
\newcommand{\R}{\mathbb{R}}
\newcommand{\C}{\mathbb{C}}
\newcommand{\F}{\mathbb{F}}
\newcommand{\A}{\mathbb{A}}
\newcommand{\PP}{\mathbb P}
\newcommand{\E}{\mathbb{E}}
\newcommand{\eps}{\varepsilon}
\newcommand{\dd}{\,\mathrm d}
\newcommand{\sref}[2]{\hyperlink{source:#1:#2}{[S#2]}}
\newcommand{\eref}[2]{\hyperlink{extra:#1:#2}{[E#2]}}
% Turn the next switch off for a shorter reading copy. The quotations preserve
% every exactTarget field, including qualifications omitted from a short summary.
\newif\ifshowcatalogue
\showcataloguetrue
\newcommand{\cataloguescope}[1]{\ifshowcatalogue
 \paragraph{Exact scope in the supplied catalogue.}
 \begin{quote}\small #1\end{quote}\fi}
\begin{document}
\setcounter{section}{152}
% ================================================================
% problemId: problem.murres-chow-kunneth-conjectures
% source releaseRank: 153; source releaseStatus: open_with_solved_subcases
% exactTarget SHA-256: 981d329e2369d80ecac265c8f104581d62a3a436182899af7698933facc122ac
\clearpage
\section{\texorpdfstring{Murre's Chow-Kunneth Conjectures}{Murre's Chow-Kunneth Conjectures}}\label{153}
\subsection{Definitions and mathematical statement}
For a smooth projective complex $d$-fold $X$, $\mathrm{CH}^j(X)_{\Q}$ is the rational vector space of codimension-$j$ cycles modulo rational equivalence. Degree-zero correspondences belong to $\mathrm{CH}^d(X\times X)_{\Q}$ and act by pullback, intersection, and pushforward. A Chow--K\"unneth decomposition consists of orthogonal idempotents $\pi_i$, summing to the diagonal, whose cohomological actions project onto $H^i(X,\Q)$. Murre's package asks for such a decomposition, vanishing of $(\pi_i)_*$ on $\mathrm{CH}^j$ outside $j\le i\le2j$, and a decomposition-independent filtration
\[
 F^r\mathrm{CH}^j=\bigcap_{i=\max(0,2j-r+1)}^{2j}\ker(\pi_i)_*,
 \qquad F^1\mathrm{CH}^j=\mathrm{CH}^j_{\mathrm{hom}}.
\]
An empty intersection is the whole group; the homologically trivial subgroup is the kernel of the Betti cycle-class map. The universal quantification over \emph{every} decomposition in B--D is retained, not replaced by existence of one specially chosen decomposition.
\par\smallskip\textit{Formulation and concept references:} \sref{153}{1}.
\cataloguescope{For every smooth projective irreducible complex variety X of dimension d, let CH\textasciicircum{}j(X)\_Q denote cycles modulo rational equivalence, with Betti Q-cohomology. A correspondence pi in CH\textasciicircum{}d(X x X)\_Q acts by pi\_*(a)=(p\_2)\_*(p\_1\textasciicircum{}*(a).pi). Murre's full A-D package asserts: (A) there exists a Chow-Kunneth decomposition, namely mutually orthogonal composition-idempotents pi\_0,...,pi\_(2d) in CH\textasciicircum{}d(X x X)\_Q summing to the diagonal and inducing projection onto H\textasciicircum{}i(X,Q) for pi\_i; (B) for EVERY such decomposition and 0<=j<=d, pi\_i acts as zero on CH\textasciicircum{}j when i<j or i>2j; (C) for EVERY such decomposition define F\textasciicircum{}r CH\textasciicircum{}j as the intersection of ker(pi\_i)\_* for max(0,2j-r+1)<=i<=2j, r>=0, with empty intersection F\textasciicircum{}0=CH\textasciicircum{}j; these filtrations are equal for ALL decompositions satisfying A; (D) for EVERY such decomposition, F\textasciicircum{}1 CH\textasciicircum{}j equals ker(CH\textasciicircum{}j(X)\_Q -> H\textasciicircum{}(2j)(X,Q)). Under A and B, F\textasciicircum{}(j+1)=0. No multiplicative or self-dual decomposition, explicit construction algorithm, functorial-extra requirement, or arbitrary-base-field enlargement is included.}
\subsection{Short English statement}
Can algebraic cycles be split into cohomological pieces so that the resulting filtration is independent of the splitting and starts precisely with the homologically trivial cycles?
\subsection{Sources}
\begin{itemize}\raggedright
\item[\textnormal{[S1]}]\hypertarget{source:153:1}{} Uwe Jannsen, Motivic Sheaves and Filtrations on Chow Groups, Proc.Symp.Pure Math.55(1994),Part1,245-302.
\url{https://epub.uni-regensburg.de/26639/1/ubr13239_ocr.pdf}.
{\small\textit{Catalogue formulation source.}}
\item[\textnormal{[E1]}]\hypertarget{extra:153:1}{} Vladimir Guletskii and
Claudio Pedrini, \emph{Finite dimensional motives and the conjectures of
Beilinson and Murre}, K-Theory 30(3) (2003), 243--263.
Theorems 8 (Kimura's nilpotence theorem) and 14 (isomorphism of the summands)
in the consulted manuscript:
\url{https://arxiv.org/html/math/0303170}.
{\small\textit{Consulted 24 September 2026; isomorphism of motives is not
identified with equality of their images in a fixed Chow group.}}
\item[\textnormal{[E2]}]\hypertarget{extra:153:2}{} Robert Laterveer and Charles
Vial, \emph{The Beauville--Voisin--Franchetta conjecture and LLSS eightfolds},
Indagationes Mathematicae 37 (2026), 250--270;
arXiv:2404.10465v1 (16 April 2024), \S1.1, Theorem 2, Proposition 1.7 and
Corollary 1.10.
\url{https://arxiv.org/html/2404.10465v1}.
Publication information:
\url{https://www.math.uni-bielefeld.de/~vial/}.
{\small\textit{Manuscript and author's publication list consulted
24 September 2026. The result on selected subrings is not used as an
injectivity theorem for the entire Chow ring.}}
\item[\textnormal{[E3]}]\hypertarget{extra:153:3}{} Charles Vial,
\emph{Remarks on motives of abelian type}, arXiv:1112.1080v3
(27 July 2015), Theorem 3 and \S2; published in Tohoku Mathematical Journal
69(2) (2017).
\url{https://arxiv.org/html/1112.1080v3}.
{\small\textit{Consulted 24 September 2026. The assumption of Murre (D)
for products of curves over a universal domain is retained.}}
% END RESEARCH ADDITION REFERENCES-153
\end{itemize}

\end{document}
